Podman ist eine Container-Engine die von Red Hat entwickelt wurde 
und als direkte Alternative zu Docker eingesetzt werden kann. 
Der Name steht für \textbf{Pod Manager} und orientiert sich 
am Konzept der Pods aus Kubernetes \cite{podman_docs}.

\subsubsection{Technische Unterschiede}

Der wesentlichste Unterschied zwischen Docker und Podman liegt 
in der Architektur. Docker benötigt einen zentralen 
\textbf{Daemon}, also einen Hintergrunddienst der mit 
Root-Rechten läuft und alle Container verwaltet. Podman 
verzichtet vollständig auf einen solchen Daemon. Stattdessen 
wird jeder Container als direkter Kindprozess des aufrufenden 
Benutzers gestartet \cite{podman_docs}.

Das hat mehrere Konsequenzen:

\begin{itemize}
    \item \textbf{Rootless Container}: Podman kann Container 
    ohne Root-Rechte starten. Das erhöht die Sicherheit, da 
    ein kompromittierter Container keine Root-Rechte auf dem 
    Host-System erlangen kann \cite{redhat_podman}.

    \item \textbf{Kein Single Point of Failure}: Da kein 
    zentraler Daemon läuft, kann ein Absturz eines Containers 
    nicht den gesamten Container-Dienst zum Absturz bringen 
    \cite{podman_docs}.

    \item \textbf{Systemd Integration}: Podman kann Container 
    direkt als Systemd-Services betreiben, was die Integration 
    in Linux-Systeme erleichtert \cite{redhat_podman}.

    \item \textbf{Pods}: Podman unterstützt nativ das Konzept 
    der Pods, also Gruppen von Containern die sich einen 
    Netzwerk-Namespace teilen, ähnlich wie bei Kubernetes 
    \cite{podman_docs}.
\end{itemize}

Ein weiterer Unterschied ist die OCI-Kompatibilität. Beide 
Tools erzeugen Images die dem OCI-Standard entsprechen und 
sind damit grundsätzlich kompatibel. Podman-Images können 
auf Docker-Registries wie Docker Hub hochgeladen und von 
Docker verwendet werden und umgekehrt \cite{redhat_podman}.

\subsubsection{Warum Docker und nicht Podman}

Für LearnHub wurde Docker gegenüber Podman aus mehreren 
Gründen bevorzugt.

\textbf{Docker Compose}: Docker Compose ist das am weitesten 
verbreitete Tool für die Orchestrierung von Multi-Container 
Anwendungen auf einem einzelnen Host. Podman bietet zwar mit 
\texttt{podman-compose} eine Alternative, diese ist jedoch 
weniger ausgereift und hat eine kleinere Community 
\cite{podman_docs}.

\textbf{Bekanntheit und Dokumentation}: Docker ist seit 2013 
der de-facto Standard für Containerisierung und hat eine 
deutlich größere Community sowie umfangreichere Dokumentation. 
Für nahezu jedes Problem gibt es bewährte Lösungen und 
Anleitungen \cite{ionos_docker}.

\textbf{Vorwissen}: Docker wurde bereits im Unterricht 
behandelt, wodurch das Entwicklungsteam auf vorhandenes 
Wissen zurückgreifen konnte. Eine Einarbeitung in Podman 
hätte zusätzliche Zeit erfordert.

\textbf{Traefik Integration}: Traefik als Reverse-Proxy 
bietet eine sehr gut dokumentierte und erprobte Integration 
mit Docker. Die automatische Service-Erkennung über den 
Docker-Socket funktioniert mit Podman zwar prinzipiell auch, 
ist aber weniger verbreitet und dokumentiert.

Podman bietet technisch interessante Vorteile, vor allem 
durch rootless Container die sicherer sind als der 
Docker-Daemon-Ansatz. Für LearnHub war Docker trotzdem die 
bessere Wahl, weil das Ökosystem ausgereifter ist, die 
Community größer ist und das Team bereits Erfahrung damit 
hatte.